\section{Stage 2: Can a Larger Model Choose Among the Candidates?}
\label{sec:cascade}
\FigureFrontier

\step{M1, the candidate generator (E14).}
On ten seeds the 12-epoch classifier scores \pmsd{0.525}{0.042} per model (+0.049 over 3 epochs, 8/10 up) and 0.598 \Stwo{} / 0.743 \Sone{} as a system (M1, Figure~\ref{fig:system}).
Three registered checks shaped the next module.
An oracle re-deciding the 30\% least confident items within the top three would reach 0.727, but a \emph{random} choice there lowers \Stwo{} to 0.587: M1's first choice carries information, so an adjudicator must be \emph{fused} with M1, not substituted for it.
Candidate sets defined by probability mass did not transfer from the slice to cross-fitted rows (coverage 0.98 vs 0.76) while rank-based sets did, so $C(x)$ is the top four.
And five seeds did not score within 0.01 of ten as predicted (0.555 vs 0.598), so M1 kept ten seeds and became the dominant cost.

\step{M3, the adjudicator (E15).}
Qwen3.8-27B-FP8, run locally without thinking, reads the label definitions, a guide to the most-confused pairs and the item, and answers with one letter among the four candidates, and $q(c \mid x)$ is the softmax over the letters' log-probabilities.
On deferred items the label is $\arg\max_{c \in C(x)} \lambda \log \bar p + (1-\lambda) \log q - \tau \log \pi_c$, with $(\lambda, \tau)$ fitted by nested CV. Other items keep M1's answer.
A screen on 300 cross-fitted training rows refuted two registered design assumptions: retrieved examples \emph{hurt} (fused accuracy 0.353 with none, 0.330 with ten), and a 304B mixture-of-experts model was worse than the local 27B (0.317).

\step{Result, and a correction.}
The first dev analysis showed no gain. Explaining it, we found a bug that shared one $\tau$ between deferred and kept items and so switched off M1's logit adjustment on the 64\% of items meant to keep M1's answer.
Corrected, deferring 30\% of items gives 0.669 \Stwo{} (0.660 over ten splits) and 0.798 \Sone{}, +0.067 \ci{+0.009}{+0.124} over M1 with the fitted parameters held fixed (Figure~\ref{fig:frontier}), with 33 items moving into the reference set and 9 out (McNemar $p = 0.0003$).
Deferring \emph{every} item is worse (0.609 over ten splits): where M1 is confident, the LLM is the weaker reader.

\step{Where the cost goes (E16).}
Timed on one GPU, an M1 seed costs 0.040 GPU-s per item and a 27B call 0.101, so the ten seeds are 92\% of the cascade's inference time and the 27B only 8\%.
Removing prompt parts lowered dev \Stwo{} in the expected order, each step within noise (full 0.660, no confusion guide 0.637, no definitions 0.628, neither 0.617), so the full prompt stayed.
With $K \le 5$ seeds M3 adds only +0.021 to +0.031 at $\delta = 0.3$, against +0.071 at $K = 10$.

\step{M4, distillation for cost (E17--E18).}
M3 relabelled every training row over its cross-fitted top four, and one student learned from $\tfrac12$gold $+ \tfrac12 q$, a weight fixed in advance \citep{hinton2015distilling}.
The student beats M1 per model (0.573 vs 0.525, 9/10 seeds up) and is far better calibrated (mean top probability 0.58 vs 0.84), yet it ranks its errors no better (AUROC 0.686 for both), against our prediction.
With M3 it missed the registered target at $\delta = 0.3$ (0.591 vs 0.614). At $\delta = 0.5$ it reaches 0.608 \Stwo{} / 0.790 \Sone{} at 0.098 GPU-s per item, above the ten-seed M1 alone at a quarter of its time, and it beats 25 of the 30 undistilled (seeds, $\delta$) points with less GPU time.
Ensembles of students add almost nothing (+0.005 from one to ten, vs +0.065 for M1 seeds): students trained towards the same targets lose the diversity seed ensembles feed on.
We froze an \emph{accuracy} point (M1 $\times$ 10 + M3, $\delta = 0.3$) and a \emph{distilled} point (one student + M3, $\delta = 0.5$).

\step{What it raised.}
The accuracy point rested on one ensemble whose M3 gain was twice that of any smaller one, on parameters fitted on 308 items, and, for the distilled point, on a $\delta$ chosen after the registered one failed.
